\chapter{Regularization}
\label{ch:regularization}

\noindent\textsc{When you have many predictors,} OLS misbehaves. With $k$ close to $n$, coefficients become wildly unstable; standard errors balloon; predictions on new data are dreadful. The model overfits: it captures noise that doesn't generalize.

The cure is regularization: deliberately bias the estimates toward zero in exchange for lower variance. The bias-variance trade-off is real, and a small bias can produce large gains in predictive accuracy. Three regularization methods—ridge, lasso, and elastic net—dominate practice.

\section{Why OLS fails with many predictors}

The OLS variance for $\hat\bbeta$ is $\sigma^2 (\bm{X}^\top \bm{X})^{-1}$. When predictors are correlated, $\bm{X}^\top \bm{X}$ is nearly singular and its inverse has huge diagonal entries. Coefficient variances explode. Predictions become noisy.

The pure cure is more data, but you usually don't have more data. The pragmatic cure is to constrain the coefficients.

\section{Ridge regression}

Add a penalty on the squared magnitude of the coefficients to the OLS objective:
\begin{equation}
  \hat\bbeta_{\text{ridge}} = \arg\min_{\bbeta}\; \|\bm{y} - \bm{X}\bbeta\|^2 + \lambda \sum_{j=1}^{k} \beta_j^2.
  \label{eq:ridge-objective}
\end{equation}
The hyperparameter $\lambda \geq 0$ controls the strength of the penalty. When $\lambda = 0$, you recover OLS. As $\lambda \to \infty$, all coefficients shrink toward zero. In between, you trade some bias for less variance.

The ridge estimator has a closed form:
\begin{equation}
  \hat\bbeta_{\text{ridge}} = (\bm{X}^\top \bm{X} + \lambda \bm{I})^{-1} \bm{X}^\top \bm{y}.
  \label{eq:ridge-estimator}
\end{equation}
The added $\lambda \bm{I}$ guarantees the matrix is invertible even when $\bm{X}^\top \bm{X}$ alone isn't. This is why ridge regression works when $k > n$ or when predictors are highly collinear: the penalty regularizes the inverse.

\paragraph{Standardize first.} The penalty depends on the units of $\beta_j$, so standardize predictors (mean 0, variance 1) before fitting. Otherwise variables measured on small scales get penalized less than variables on large scales, for no good reason.

\section{Lasso}

Replace the squared-coefficient penalty with absolute values:
\begin{equation}
  \hat\bbeta_{\text{lasso}} = \arg\min_{\bbeta}\; \|\bm{y} - \bm{X}\bbeta\|^2 + \lambda \sum_{j=1}^{k} |\beta_j|.
  \label{eq:lasso-objective}
\end{equation}
The $\ell_1$ penalty does something the $\ell_2$ penalty does not: it pushes some coefficients exactly to zero. The lasso simultaneously shrinks coefficients and performs variable selection.

Geometrically, the $\ell_1$ ball has corners along the axes. As you increase the penalty, the optimal solution gets pulled toward these corners, where some coefficients are exactly zero.

There's no closed form. Lasso is solved by coordinate descent or LARS. R's \texttt{glmnet} package handles this fast even for thousands of predictors.

\begin{figure}[ht]
\centering
\includegraphics[width=\linewidth]{fig_regularization.pdf}
\caption{Coefficient paths as the penalty $\lambda$ varies. Left (ridge): coefficients shrink smoothly toward zero but never reach it. Right (lasso): coefficients hit zero at finite $\lambda$, performing variable selection. Solid lines: truly nonzero coefficients. Dashed lines: truly zero coefficients (noise). Lasso correctly drives the noise coefficients to exactly zero.}
\label{fig:regularization}
\end{figure}



\section{Elastic net}

Combine both penalties:
\begin{equation}
  \hat\bbeta_{\text{enet}} = \arg\min_{\bbeta}\; \|\bm{y} - \bm{X}\bbeta\|^2 + \lambda_1 \sum_j |\beta_j| + \lambda_2 \sum_j \beta_j^2.
\end{equation}
Combines lasso's variable selection with ridge's stability when predictors are highly correlated. In practice, parameterized as a single $\lambda$ and a mixing parameter $\alpha \in [0, 1]$ (where $\alpha = 1$ is pure lasso and $\alpha = 0$ is pure ridge).

\section{Choosing $\lambda$}

The penalty strength is a hyperparameter. Too small, you overfit; too large, you underfit. Choose by cross-validation:

\begin{enumerate}[noitemsep]
\item Split the data into $K$ folds (typically $K = 5$ or $10$).
\item For each candidate $\lambda$ value: fit the model on $K-1$ folds, predict on the held-out fold, record the test error.
\item Average across folds. The $\lambda$ with the lowest cross-validated error wins.
\end{enumerate}

The R command \texttt{cv.glmnet()} does this automatically. The result is the value of $\lambda$ that gives the best out-of-sample prediction.

\section{The bias-variance trade-off}

Why do these methods help? The expected prediction error of an estimator decomposes:
\begin{equation}
  \E\bigl[(y - \hat y)^2\bigr] = \text{Bias}^2 + \text{Variance} + \text{Irreducible noise}.
  \label{eq:bias-variance}
\end{equation}
OLS is unbiased ($\text{Bias} = 0$) but high variance. Ridge introduces bias to reduce variance. The total prediction error can be lower for the biased estimator, especially when predictors are correlated or $k$ is large relative to $n$.

This is one of the deepest insights in modern statistics. Unbiasedness is not the goal. \emph{Low expected prediction error} is the goal, and a little bias is often the price of getting there.

\begin{keyinsight}
OLS is unbiased but can have huge variance when predictors are correlated or numerous.

Regularization adds a penalty to the OLS objective, biasing estimates toward zero in exchange for much lower variance. Ridge ($\ell_2$ penalty) shrinks all coefficients smoothly. Lasso ($\ell_1$) shrinks and selects: it sets some coefficients exactly to zero. Elastic net combines both.

Choose the penalty strength $\lambda$ by cross-validation. The bias-variance trade-off means a slightly biased estimator can have lower expected prediction error than the unbiased OLS.
\end{keyinsight}

\begin{codelab}
\begin{lstlisting}[language=Rlang]
# install.packages("glmnet")
library(glmnet)

# Simulate a high-dimensional problem
set.seed(2024)
n <- 100; p <- 50
X <- matrix(rnorm(n * p), n, p)
beta_true <- c(rep(2, 5), rep(0, 45))    # only 5 nonzero
y <- X %*% beta_true + rnorm(n)

# Ridge
fit_ridge <- glmnet(X, y, alpha = 0)        # alpha = 0 is ridge
plot(fit_ridge, xvar = "lambda")            # coefficient paths

# Lasso
fit_lasso <- glmnet(X, y, alpha = 1)        # alpha = 1 is lasso
plot(fit_lasso, xvar = "lambda")

# Cross-validation for lambda
cv_lasso <- cv.glmnet(X, y, alpha = 1)
plot(cv_lasso)
cv_lasso$lambda.min                          # best lambda
cv_lasso$lambda.1se                          # one SE rule

# Coefficients at chosen lambda
coef(cv_lasso, s = "lambda.min")             # nonzero ones survive
\end{lstlisting}
\end{codelab}

\section{Practice problems}

\begin{enumerate}
  \item Show that the ridge estimator equals $(\bm{X}^\top \bm{X} + \lambda \bm{I})^{-1}\bm{X}^\top \bm{y}$. Take the derivative of the ridge objective and set it to zero.
  \item For a single-predictor model, show that $\hat\beta_{\text{ridge}} = \hat\beta_{\text{OLS}} \cdot \frac{S_{xx}}{S_{xx} + \lambda}$. Notice that ridge always shrinks toward zero.
  \item Why does lasso produce exactly-zero coefficients but ridge doesn't? Sketch the constraint sets ($\sum |\beta_j| \leq t$ vs.\ $\sum \beta_j^2 \leq t$) in two dimensions.
  \item Generate a high-dimensional dataset where only a few predictors matter. Fit OLS, ridge, and lasso. Compare which method recovers the truth and which overfits.
  \item Why must you standardize predictors before applying regularization? Construct an example where failing to standardize gives misleading coefficient comparisons.
  \item Use \texttt{cv.glmnet} to choose $\lambda$ for a real or simulated dataset. Compare the model selected by lambda.min vs.\ lambda.1se.
\end{enumerate}

\begin{reflect}
$\triangleright$~State the bias-variance trade-off in your own words. Why is unbiasedness not always the goal?

\smallskip
$\triangleright$~Lasso does variable selection. Ridge does not. What property of the $\ell_1$ penalty makes this happen?

\smallskip
$\triangleright$~Cross-validation chooses the penalty strength. Why does it work better than just looking at the training error?
\end{reflect}
